\section{Introduction}
\label{sec:introduction}

Speech emotion recognition (SER) models are usually trained on one corpus and
deployed on another. Differences in speakers, elicitation procedures, recording
conditions and label definitions between corpora degrade cross-corpus
performance \cite{schuller2010cross}. Unsupervised domain adaptation addresses
this degradation by aligning the feature distributions of the source (training)
and target (deployment) corpora, for example by matching their means,
covariances or kernel mean embeddings
\cite{sun2016coral,gretton2012kernel,long2015dan}. Every alignment method has
design choices and hyperparameters, and in deployment these must be chosen
without target labels. The standard practice is to select them on held-out
\emph{source} data. Source-side validation is procedurally safe because it uses
no target label and therefore cannot leak target information into model
selection.

Procedural safety, however, does not imply that the criterion is informative.
Unsupervised model selection under domain shift is known to be difficult, and
several studies propose or evaluate selection criteria that use signals beyond
the source-validation score
\cite{you2019selection,ericsson2023practices,hu2024selection}. Matching marginal
feature distributions while keeping the source error low is also known to be
insufficient for successful adaptation \cite{zhao2019invariant}. Two questions
nevertheless remain open for cross-corpus SER. First, to our knowledge, it
has not been characterised exactly when source-side validation carries no
information about an adaptation choice. Second, a comparison of alignment methods ordered by the
moments they match, such as the experimental grid analysed in this work,
implicitly treats the methods as otherwise-identical pipelines. An affine map,
however, also changes the coordinates in which the classifier is fitted and in
which a distribution discrepancy is measured.

This work therefore addresses two problems: identifying when a leakage-free
source-validation criterion cannot distinguish between affine feature-alignment
choices, and determining what a broader comparison of affine alignment methods,
together with its discrepancy diagnostics, can identify. We combine an
analytical result with a retrospective audit of completed experiments. We first
show that, for a classifier whose fit depends on the source features only
through a stationary kernel matrix, such as a support vector machine (SVM) with
a radial basis function (RBF) kernel, translating all source training and
validation features by a common offset leaves every source-validation
prediction unchanged. The target predictions, in contrast, can change.
Source-side validation is therefore exactly constant across adaptation choices
with different target behaviour, and a larger labelled validation set cannot
resolve this tie.

We test this prediction on a completed grid of 4,986 cross-corpus runs between
RAVDESS \cite{livingstone2018ryerson} and CREMA-D \cite{cao2014crema}, covering
three self-supervised speech encoders, six alignment conditions and five
classifier families. The configuration and the result ledger of this grid were
frozen, that is, fixed under a version tag and protected by a checksum, before
the audit was designed. The audit is read-only. It matches unaligned and
mean-shifted runs and compares their stored validation scores and selected
hyperparameters without rerunning any experiment. Finally, we factor a general
affine adaptation into relative transport between the domains and a common
change of coordinates. The common part changes the classifier's effective
regularisation and the geometry in which a discrepancy is measured.
Figure~\ref{fig:pipeline} summarises the experimental pipeline and the audit.

\begin{figure*}[!t]
  \centering
  \includegraphics[width=\textwidth]{pipeline.pdf}
  \caption{Overview of the experimental pipeline (top) and the retrospective
    read-only audit (bottom).}
  \label{fig:pipeline}
\end{figure*}

The audit confirms the prediction exactly. In every matched RBF-SVM cell in
both transfer directions, the stored validation score and the selected
hyperparameters are identical with and without mean shift, although mean shift
changes target macro-F1 in every cell. Logistic regression, a linear SVM and a
multilayer perceptron (MLP), whose fitting procedures fall outside the stated
assumptions, do not generally retain exact equality. At the level of the complete
protocol, source-only selection chooses the unaligned condition on two of five
seeds in the RAVDESS-to-CREMA-D direction. On one of these seeds, the selected
configuration scores below the chance baseline.

Cross-corpus degradation, the insufficiency of marginal invariance and the
difficulty of unsupervised model selection are established
\cite{schuller2010cross,zhao2019invariant,you2019selection,ericsson2023practices,hu2024selection}.
The affine reparameterisation used in this work is also elementary and already
implicit in the CORAL account \cite{sun2016coral}. The novelty of this work lies
in identifying an \emph{exact} symmetry of a leakage-free selection criterion
under stated assumptions, confirming it in a controlled SER experiment, and
bounding how far the result extends. Our contributions are as follows:
\begin{enumerate}
 \item We establish a source-translation invariance of source-side validation
   for stationary-kernel classifiers. We state its fitting assumptions, its
   consequence for model selection, and the numerical boundary at which it
   stops holding (Section~\ref{sec:translation-theory}).
 \item We provide an executable, read-only audit of matched experimental runs.
   The audit tests the predicted equality without rounding and without
   target-based cell selection, and it covers every classifier family with
   balanced matched cells, so that the boundary of the result is measured
   rather than asserted
   (Sections~\ref{sec:results-translation} and~\ref{sec:results-boundary}).
 \item We factor affine adaptation into relative transport, classifier
   coordinates and measurement geometry. The factorisation states what a
   broader affine comparison and its discrepancy diagnostics can and cannot
   identify
   (Sections~\ref{sec:affine-factorisation} and~\ref{sec:results-frames}).
\end{enumerate}

This work does not propose a new recogniser or an unsupervised estimator of
target risk. Instead, it characterises a blind spot in a criterion that is
widely used precisely because it is safe. The remainder of the paper is
organised as follows. Section~\ref{sec:related} reviews related work.
Section~\ref{sec:selection-theory} presents the analysis, and
Section~\ref{sec:method} describes the experimental design.
Section~\ref{sec:results} reports the audit and the surrounding comparisons.
Section~\ref{sec:discussion} discusses the implications and limitations, and
Section~\ref{sec:conclusion} concludes.
